\documentclass[10pt,a4paper]{amsart}
\usepackage{amsmath,amssymb,amsthm,graphicx}
% Retained minimal preamble for the Stacks Project chapter flat.tex (More on Flatness).
% Verbatim excerpt of the chapter preamble: Xy-pic setup (native preamble lines 14-18) and the native
% theorem declarations with the mathematical macros (native preamble lines 154-217). Copied mechanically;
% nothing added, removed or redefined. The wrapper class and the amsmath/amssymb/amsthm/graphicx base
% packages are supplied by the assembler, so the class/externaldocument part of the native preamble is omitted.
% For commutative diagrams we use Xy-pic
\usepackage[all]{xy}

% We use 2cell for 2-commutative diagrams.
\xyoption{2cell}
% Theorem environments.
%
\theoremstyle{plain}
\newtheorem{theorem}[subsection]{Theorem}
\newtheorem{proposition}[subsection]{Proposition}
\newtheorem{lemma}[subsection]{Lemma}

\theoremstyle{definition}
\newtheorem{definition}[subsection]{Definition}
\newtheorem{example}[subsection]{Example}
\newtheorem{exercise}[subsection]{Exercise}
\newtheorem{situation}[subsection]{Situation}

\theoremstyle{remark}
\newtheorem{remark}[subsection]{Remark}
\newtheorem{remarks}[subsection]{Remarks}

\numberwithin{equation}{subsection}

% Macros
%
\def\lim{\mathop{\mathrm{lim}}\nolimits}
\def\colim{\mathop{\mathrm{colim}}\nolimits}
\def\Spec{\mathop{\mathrm{Spec}}}
\def\Hom{\mathop{\mathrm{Hom}}\nolimits}
\def\Ext{\mathop{\mathrm{Ext}}\nolimits}
\def\SheafHom{\mathop{\mathcal{H}\!\mathit{om}}\nolimits}
\def\SheafExt{\mathop{\mathcal{E}\!\mathit{xt}}\nolimits}
\def\Sch{\mathit{Sch}}
\def\Mor{\mathop{\mathrm{Mor}}\nolimits}
\def\Ob{\mathop{\mathrm{Ob}}\nolimits}
\def\Sh{\mathop{\mathit{Sh}}\nolimits}
\def\NL{\mathop{N\!L}\nolimits}
\def\CH{\mathop{\mathrm{CH}}\nolimits}
\def\proetale{{pro\text{-}\acute{e}tale}}
\def\etale{{\acute{e}tale}}
\def\QCoh{\mathit{QCoh}}
\def\Ker{\mathop{\mathrm{Ker}}}
\def\Im{\mathop{\mathrm{Im}}}
\def\Coker{\mathop{\mathrm{Coker}}}
\def\Coim{\mathop{\mathrm{Coim}}}

% Boxtimes
%
\DeclareMathSymbol{\boxtimes}{\mathbin}{AMSa}{"02}

%
% Macros for moduli stacks/spaces
%
\def\QCohstack{\mathcal{QC}\!\mathit{oh}}
\def\Cohstack{\mathcal{C}\!\mathit{oh}}
\def\Spacesstack{\mathcal{S}\!\mathit{paces}}
\def\Quotfunctor{\mathrm{Quot}}
\def\Hilbfunctor{\mathrm{Hilb}}
\def\Curvesstack{\mathcal{C}\!\mathit{urves}}
\def\Polarizedstack{\mathcal{P}\!\mathit{olarized}}
\def\Complexesstack{\mathcal{C}\!\mathit{omplexes}}
% \Pic is the operator that assigns to X its picard group, usage \Pic(X)
% \Picardstack_{X/B} denotes the Picard stack of X over B
% \Picardfunctor_{X/B} denotes the Picard functor of X over B
\def\Pic{\mathop{\mathrm{Pic}}\nolimits}
\def\Picardstack{\mathcal{P}\!\mathit{ic}}
\def\Picardfunctor{\mathrm{Pic}}
\def\Deformationcategory{\mathcal{D}\!\mathit{ef}}

% --- appended 2026-10-12 (Chief): the project own definitions of the reference/slogan/history
% display environments, copied mechanically from the shared Stacks preamble lines 8-12
% (\usepackage{verbatim} plus the three \newenvironment lines, which the chapter preamble keeps
% verbatim from the project preamble). Nothing added, removed or redefined: these environments
% carry no mathematical content and the project itself typesets them as comments. Without these
% three lines a counted statement whose own text uses them fails with Environment slogan undefined.
\usepackage{verbatim}
\newenvironment{reference}{\comment}{\endcomment}
\newenvironment{slogan}{\comment}{\endcomment}
\newenvironment{history}{\comment}{\endcomment}
\usepackage{cleveref}
\crefname{originallemma}{original lemma}{original lemmas}
\Crefname{originallemma}{Original lemma}{Original lemmas}
\makeatletter
\newlabel{algebra-lemma-finite-flat-local}{{00NZ}{source}{}{sourcepointer-1}{}}
\newlabel{algebra-lemma-finite-flat-local@cref}{{[originallemma][0][]00NZ}{[1][0][]source}{}{}{}}
\makeatother
\crefname{originallemma}{original lemma}{original lemmas}
\Crefname{originallemma}{Original lemma}{Original lemmas}
\makeatletter
\newlabel{divisors-lemma-weakly-ass-pullback}{{05F0}{source}{}{sourcepointer-2}{}}
\newlabel{divisors-lemma-weakly-ass-pullback@cref}{{[originallemma][0][]05F0}{[1][0][]source}{}{}{}}
\makeatother
\crefname{originallemma}{original lemma}{original lemmas}
\Crefname{originallemma}{Original lemma}{Original lemmas}
\makeatletter
\newlabel{divisors-lemma-weakly-associated-finite}{{05EZ}{source}{}{sourcepointer-3}{}}
\newlabel{divisors-lemma-weakly-associated-finite@cref}{{[originallemma][0][]05EZ}{[1][0][]source}{}{}{}}
\makeatother
\title{Weakly associated points of a finite flat pullback}
\author{The Stacks Project Authors}
\date{stacks-project:flat.tex:local-snapshot-20261011; GNU FDL 1.2 or later}
\begin{document}
\maketitle
\noindent\textit{Editorial note.} All delivered mathematical text is the author's own, copied from the pinned Stacks Project chapter file flat.tex; no statement or proof step was rewritten, reordered or removed.\par\medskip
\noindent\textit{Editorial note.} Preamble provenance: the chapter's own retained minimal preamble (sources/stacks-flat/retained-preamble.tex, lines 1-85) is a verbatim excerpt of the native Xy-pic setup and of the native theorem declarations and mathematical macros, followed by the project's own reference/slogan/history display-environment definitions copied mechanically from the shared preamble; the wrapper class and the amsmath/amssymb/amsthm/graphicx base packages are supplied by the assembler, and extra\_packages is left empty because everything needed is already in the retained preamble. No macro was invented or redefined.\par\medskip
\noindent\textit{Editorial note.} Same-chapter pointers are located at their \textbackslash{}label lines inside the pinned chapter file; cross-chapter pointers are located in the project's own tag list staged in the snapshot (sources/stacks-chow/tags.txt), at the verbatim lines carrying their canonical tags.\par\medskip
\noindent\textit{Allowed standard prerequisites (editorial).}\begin{itemize}
\item Weakly associated primes under a flat pullback. Divisors, label divisors-lemma-weakly-ass-pullback, invoked at proof line 325. Gives the easy implication: a weakly associated prime downstairs pulls back to one upstairs because the morphism is flat.
\item A finite flat algebra over a local ring is free, and weakly associated points under a finite morphism. Algebra, label algebra-lemma-finite-flat-local (proof line 336); Divisors, label divisors-lemma-weakly-associated-finite (proof line 348). The first gives that the localised algebra is a free module of some rank n at least one, so that the pushforward of the pulled back module has the same weakly associated primes as the original downstairs, and the second pushes a weakly associated point through the finite morphism.
\end{itemize}\medskip
\ifdefined\hypertarget\hypertarget{sourcepointer-1}{}\fi
\noindent\textit{Source reference.} Original lemma 00NZ: sources/stacks-chow/tags.txt, source lines 877--877. Navigation pointer for \textbackslash{}ref\{algebra-lemma-finite-flat-local\} used inside the counted span (line 336); the argument that consumes it is retained in the span.\par\medskip
\ifdefined\hypertarget\hypertarget{sourcepointer-2}{}\fi
\noindent\textit{Source reference.} Original lemma 05F0: sources/stacks-chow/tags.txt, source lines 7009--7009. Navigation pointer for \textbackslash{}ref\{divisors-lemma-weakly-ass-pullback\} used inside the counted span (line 325); the argument that consumes it is retained in the span.\par\medskip
\ifdefined\hypertarget\hypertarget{sourcepointer-3}{}\fi
\noindent\textit{Source reference.} Original lemma 05EZ: sources/stacks-chow/tags.txt, source lines 7008--7008. Navigation pointer for \textbackslash{}ref\{divisors-lemma-weakly-associated-finite\} used inside the counted span (line 348); the argument that consumes it is retained in the span.\par\medskip
\begin{lemma}
\label{lemma-finite-flat-weak-assassin-up-down}
Let $g : T \to S$ be a finite flat morphism of schemes.
Let $\mathcal{G}$ be a quasi-coherent $\mathcal{O}_S$-module.
Let $t \in T$ be a point with image $s \in S$. Then
$$
t \in \text{WeakAss}(g^*\mathcal{G})
\Leftrightarrow
s \in \text{WeakAss}(\mathcal{G})
$$
\end{lemma}

\begin{proof}
The implication ``$\Leftarrow$'' follows immediately from
Divisors, Lemma \ref{divisors-lemma-weakly-ass-pullback}.
Assume $t \in \text{WeakAss}(g^*\mathcal{G})$.
Let $\Spec(A) \subset S$ be an affine open neighbourhood of $s$.
Let $\mathcal{G}$ be the quasi-coherent sheaf associated to the $A$-module $M$.
Let $\mathfrak p \subset A$ be the prime ideal corresponding to $s$.
As $g$ is finite flat we have $g^{-1}(\Spec(A)) = \Spec(B)$
for some finite flat $A$-algebra $B$. Note that
$g^*\mathcal{G}$ is the quasi-coherent $\mathcal{O}_{\Spec(B)}$-module
associated to the $B$-module $M \otimes_A B$ and $g_*g^*\mathcal{G}$ is the
quasi-coherent $\mathcal{O}_{\Spec(A)}$-module associated to the
$A$-module $M \otimes_A B$. By
Algebra, Lemma \ref{algebra-lemma-finite-flat-local}
we have $B_{\mathfrak p} \cong A_{\mathfrak p}^{\oplus n}$
for some integer $n \geq 0$. Note that $n \geq 1$ as we assumed there
exists at least one point of $T$ lying over $s$. Hence we see by
looking at stalks that
$$
s \in \text{WeakAss}(\mathcal{G})
\Leftrightarrow
s \in \text{WeakAss}(g_*g^*\mathcal{G})
$$
Now the assumption that $t \in \text{WeakAss}(g^*\mathcal{G})$
implies that $s \in \text{WeakAss}(g_*g^*\mathcal{G})$ by
Divisors, Lemma \ref{divisors-lemma-weakly-associated-finite}
and hence by the above $s \in  \text{WeakAss}(\mathcal{G})$.
\end{proof}

\end{document}
